\documentclass[11pt,a4paper]{article}
\usepackage[top=2cm,bottom=2cm,left=2.5cm,right=2.5cm]{geometry}
\usepackage{booktabs}
\usepackage{enumitem}
\usepackage{parskip}
\usepackage{titlesec}
\usepackage{hyperref}

\setlist{nosep, leftmargin=1.5em, topsep=2pt}
\titlespacing*{\section}{0pt}{6pt}{3pt}
\titlespacing*{\subsection}{0pt}{5pt}{2pt}
\setlength{\parskip}{4pt}
\pagestyle{empty}

\title{\vspace{-2.5em}\textbf{ML Project Plan: Restaurant Preparation Time Estimation}\vspace{-1.5em}}
\author{}
\date{}

\begin{document}
\maketitle
\vspace{-1em}

\section{Introduction}

Deliveroo's delivery experience depends on precise coordination between food preparation and rider dispatch. Currently, a restaurant's preparation time is estimated as the average of all its orders over the last month. This ignores rich structure in the data---a 20-item Friday-evening pizza order takes far longer than a single weekday lunch salad, yet both receive the same estimate. \textbf{Overestimates} cause riders to wait idle (increased cost, reduced fleet efficiency); \textbf{underestimates} cause late arrivals (cold food, poor CX, lower retention). A better model improves all three sides of the marketplace.

\section{Goals / Non-goals}

\textbf{Goals:}
\begin{itemize}
\item Build a predictive model for \texttt{prep\_time\_seconds} that significantly outperforms the per-restaurant monthly average.
\item Capture key drivers: order complexity, restaurant characteristics, food type, temporal patterns, kitchen busyness.
\item Produce well-calibrated, interpretable estimates restaurant partners can understand and trust.
\end{itemize}

\textbf{Non-goals:} Predicting rider travel time (separate model). Real-time mid-preparation updates (future extension). Prescribing operational changes to restaurants.

\section{Proposed Approach}

\subsection{Data and Feature Engineering}

\textbf{Target variable:} \texttt{prep\_time\_seconds} (continuous, non-negative, right-skewed---likely log-normal or gamma).

\vspace{-0.5em}
\begin{table}[h]
\centering\small
\begin{tabular}{@{}p{2.8cm}p{4.2cm}p{7.5cm}@{}}
\toprule
\textbf{Feature Group} & \textbf{Features} & \textbf{Rationale} \\
\midrule
Order complexity & \texttt{number\_of\_items}, \texttt{order\_value\_gbp}, value-per-item & More items = longer prep. Value proxies for dish complexity. \\
Food type & \texttt{type\_of\_food} (target-encoded) & Pizza vs.\ sushi vs.\ salad have fundamentally different prep processes. \\
Restaurant ID & \texttt{restaurant\_id} (target-encoded) & Captures unobserved factors: kitchen size, staffing, equipment. \\
Geography & \texttt{country}, \texttt{city} & Cultural and operational differences in kitchen norms. \\
Temporal & Hour (shift bins), day of week, month, country-specific holidays & Non-linear demand; country-specific calendars (Ramadan, bank holidays, festivals). \\
Interactions & food $\times$ restaurant, calendar $\times$ country & Turkish restaurant's shawarma prep $\neq$ its rare pizza. Holidays vary by country. \\
Busyness & Concurrent orders at \texttt{order\_acknowledged\_at} & Key latent driver: count overlapping [acknowledged, ready] windows. \\
Historical stats & Trailing-window median, volume, variance & Proxies for kitchen capacity; jumps signal operational changes. \\
\bottomrule
\end{tabular}
\end{table}
\vspace{-0.5em}

\textbf{Processing:} (1)~Remove outliers ($<$30s or $>$2hr). (2)~Parse timestamps; compute concurrent orders. (3)~Standardise \texttt{type\_of\_food}, consolidate duplicates, target-encode. (4)~Temporal train/val/test split to avoid leakage.

\subsection{Modelling Approach}

An \textbf{iterative approach} with increasing complexity; each stage evaluated before proceeding:

\textbf{Stage 1---Enhanced baselines:} Per-restaurant average segmented by time-of-day and item-count buckets. Linear regression on engineered features (log-transformed target). Establishes feature importance.

\textbf{Stage 2---Gradient-boosted trees (primary):} LightGBM/XGBoost on the full feature set. Handles non-linearity, interactions, and mixed types naturally. \textit{Quantile regression} variant (10th/50th/90th percentiles) produces calibrated intervals---critical because over- and under-estimates have \textbf{asymmetric costs}. The optimal prediction may be a specific quantile (e.g.\ 60th percentile) rather than the mean, tuned via the dispatch simulator.

\textbf{Stage 3---Mixed-effects model:} Restaurant as random effect, order features as fixed effects. Handles hierarchical structure and provides calibrated uncertainty. Valuable for \textbf{cold-start} restaurants: kitchen parameters (staff, stations) initialise informative priors.

\textbf{Stage 4 (exploratory)---Survival analysis:} Prep time as time-to-event. Enables \textbf{dynamic prediction}: given $t$\,min already elapsed, what is the expected remaining time? Operationally valuable for real-time dispatch. Parametric (Weibull) or neural (mixture density network) variants.

\textbf{Iteration rule:} If Stage~2 achieves $<$15\% MAE improvement over Stage~1, invest in feature engineering before adding model complexity.

\subsection{Model Explainability}

\textbf{Global:} SHAP summary plots identify dominant features; if \texttt{restaurant\_id} dominates, the model is memorising rather than generalising. \textbf{Local:} Per-order SHAP waterfall plots power restaurant-facing insights: \textit{``+8\,min from 6 concurrent orders, +4\,min Friday rush, $-$2\,min your fast kitchen.''} \textbf{Drift detection:} Significant shifts in feature importance between retraining cycles flag distribution changes.

\subsection{Model Evaluation}

\textbf{Statistical metrics:} Primary: \textbf{MAE} (interpretable in seconds, robust to outliers). Secondary: RMSE (penalises large errors); quantile calibration. Evaluate per-segment (country, food type, restaurant tier, time of day) to ensure uniform improvement.

\textbf{Business metrics:} (1)~\textbf{Rider wait time} (primary KPI): simulate dispatch with new vs.\ baseline model. (2)~Early arrival rate (\% $>$5\,min before food ready). (3)~Late arrival rate (\% $>$5\,min after). (4)~Food-ready-to-pickup latency.

\textbf{Measuring impact:} \textit{Offline}---replay historical orders through dispatch simulator. \textit{Online}---A/B test on random restaurant/city subset, measuring wait time, ratings, and rider efficiency (deliveries/hr) with cluster-corrected significance tests.

\textbf{Monitoring:} Retrain weekly on rolling window. Flag restaurants whose MAE exceeds $2\times$ population median for two consecutive weeks.

\end{document}
